\documentclass[aspectratio=169]{beamer}
\usetheme{Madrid}
\usecolortheme{whale}
\usepackage[utf8]{inputenc}
\usepackage[T1]{fontenc}
\usepackage{lmodern}
\usepackage[spanish]{babel}
\usepackage{booktabs}
\usepackage{graphicx}
\usepackage{tikz}
\usetikzlibrary{arrows.meta,positioning}
\graphicspath{{images/}}
\setbeamertemplate{navigation symbols}{}
\pdfstringdefDisableCommands{\def\\{ }}
\title[CNN \& TL Challenges]{CNN \& Transfer Learning Challenges}
\subtitle{Experimentos, ajustes y resultados en clasificación de imágenes}
\author[Gómez, Miralles, Sánchez]{Gómez Hernández, Marco \\
Miralles Retuerto, Pablo \\ Sánchez Carmona, Germán}
\institute[UPM]{Universidad Politécnica de Madrid\\
Máster Universitario en Inteligencia Artificial\\Visión por Computador}
\date{8 de noviembre de 2026}

\newcommand{\idea}[1]{\begin{center}\large\color{blue!65!black}#1\end{center}}
\newcommand{\grafica}[2]{\begin{center}\includegraphics[width=\linewidth,height=#2,keepaspectratio]{#1}\end{center}}
\begin{document}
% 1: 15 s
\begin{frame}\titlepage\end{frame}

\section{CNN Challenge}
% 2: 55 s
\begin{frame}{CNN Challenge · Punto de partida: cuatro populares}
\centering\small
\textbf{Desde cero} · 128 px · batch 32 · AdamW · LR $10^{-3}$ · WD $10^{-4}$
\grafica{populares_curvas.pdf}{4.05cm}
\begin{columns}[T,onlytextwidth]
\begin{column}{.48\textwidth}\centering
\textbf{Seguimos con ResNet18}\\\small 78,03\,\% · entrenamiento más rápido
\end{column}
\begin{column}{.48\textwidth}\centering
\textbf{Y con EfficientNetB0}\\\small 77,97\,\% · menor sobreajuste inicial
\end{column}
\end{columns}
\end{frame}

% 3: 60 s
\begin{frame}{CNN Challenge · Ajustes de las populares}
\grafica{ajustes_curvas.pdf}{3.9cm}
\begin{columns}[T,onlytextwidth]
\begin{column}{.32\textwidth}\centering
\textbf{Exp02 · No mejoran}\\\small
ResNet: WD $10^{-3}$, 78,03\,\%\\
B0: LR $3\cdot10^{-4}$, 76,75\,\%
\end{column}
\begin{column}{.32\textwidth}\centering
\textbf{Exp03 · Smoothing 0.1}\\\small
ResNet: \textbf{79,73\,\%}\\
B0: \textbf{79,09\,\%}
\end{column}
\begin{column}{.32\textwidth}\centering
\textbf{Exp04 · ResNet a 192 px}\\\small
Conserva smoothing 0.1\\
\textbf{80,11\,\%}: mejor popular
\end{column}
\end{columns}
\end{frame}

% 4: 45 s
\begin{frame}{CNN Challenge · Pasamos a una arquitectura propia}
\begin{center}
\begin{tikzpicture}[node distance=.25cm,box/.style={draw=blue!65!black,fill=blue!6,rounded corners,align=center,font=\small,minimum height=1cm},arr/.style={-{Latex[length=2mm]},thick,blue!65!black}]
\node[box,minimum width=1.3cm] (in) {RGB\\$128^2$};
\node[box,right=of in,minimum width=1.5cm] (stem) {Stem\\2 conv};
\node[box,right=of stem,minimum width=1.65cm] (s1) {2 bloques SE\\32 canales};
\node[box,right=of s1,minimum width=1.65cm] (s2) {2 bloques SE\\64 canales};
\node[box,below=.45cm of s2,minimum width=1.65cm] (s3) {2 bloques SE\\128 canales};
\node[box,left=of s3,minimum width=1.65cm] (s4) {2 bloques SE\\256 canales};
\node[box,left=of s4,minimum width=1.5cm] (gap) {GAP};
\node[box,left=of gap,minimum width=1.3cm] (out) {Dense13\\Softmax};
\draw[arr] (in)--(stem);\draw[arr] (stem)--(s1);\draw[arr] (s1)--(s2);\draw[arr] (s2)--(s3);\draw[arr] (s3)--(s4);\draw[arr] (s4)--(gap);\draw[arr] (gap)--(out);
\end{tikzpicture}
\end{center}
\begin{columns}[c,onlytextwidth]
\begin{column}{.48\textwidth}\centering
{\LARGE\color{blue!65!black}2,86 M}\\\small parámetros
\end{column}
\begin{column}{.48\textwidth}\centering
{\LARGE\color{blue!65!black}79,15\,\%}\\\small accuracy inicial · custom exp01
\end{column}
\end{columns}
\vspace{.2cm}\centering\small Mismos hiperparámetros iniciales de las populares · sin preentrenamiento
\end{frame}

% 5: 50 s
\begin{frame}{CNN Challenge · Custom: qué regularización conservar}
\grafica{custom_inicio.pdf}{4.1cm}
\begin{columns}[T,onlytextwidth]
\begin{column}{.32\textwidth}\centering
\textbf{Exp01 · Referencia}\\79,15\,\%
\end{column}
\begin{column}{.32\textwidth}\centering
\textbf{Exp02 · Smoothing 0.1}\\79,04\,\% · se descarta
\end{column}
\begin{column}{.32\textwidth}\centering
\textbf{Exp03 · MixUp 0.2}\\\textbf{80,69\,\%} · se conserva
\end{column}
\end{columns}
\end{frame}

% 6: 45 s
\begin{frame}{CNN Challenge · Custom: las dos últimas pruebas}
\grafica{custom_final.pdf}{4.05cm}
\begin{columns}[T,onlytextwidth]
\begin{column}{.32\textwidth}\centering
\textbf{Exp03 · Referencia}\\\small MixUp 0.2 · 128 px\\\textbf{80,69\,\%}
\end{column}
\begin{column}{.32\textwidth}\centering
\textbf{Exp04 · MixUp 0.4}\\\small Mantiene 128 px\\80,16\,\% · no mejora
\end{column}
\begin{column}{.32\textwidth}\centering
\textbf{Exp05 · 192 px}\\\small Vuelve a MixUp 0.2\\79,47\,\% · no mejora
\end{column}
\end{columns}
\end{frame}

\section{TL Challenge}
% 7: 55 s
\begin{frame}{TL Challenge · Punto de partida: EfficientNetB0}
\begin{columns}[T,onlytextwidth]
\begin{column}{.48\textwidth}
\begin{block}{¿Por qué EfficientNetB0?}
\begin{itemize}
\item Buenos resultados en CNN.
\item Pesos ImageNet en Keras.
\item ResNet18 no está en Keras Applications.
\end{itemize}
\end{block}
\end{column}
\begin{column}{.48\textwidth}
\begin{block}{Exp01 · Extractor congelado}
128 px · GAP + Dense13\\[6pt]
Solo se entrena la cabeza\\[6pt]
20 épocas · LR $10^{-3}$\\[6pt]
\textbf{70,72\,\%} de accuracy
\end{block}
\end{column}
\end{columns}
\vspace{.3cm}
\idea{Referencia B0 desde cero: \textbf{77,97\,\%}}
\end{frame}

% 8: 65 s
\begin{frame}{TL Challenge · De los últimos bloques a toda la red}
\centering\small 128 px · LR $10^{-5}$ · WD $10^{-4}$ · batch 32 · BN congeladas
\grafica{tl_adaptacion.pdf}{3.95cm}
\begin{columns}[T,onlytextwidth]
\begin{column}{.32\textwidth}\centering
\textbf{Exp02 · FT parcial}\\\small Últimos bloques + cabeza\\\textbf{75,20\,\%}
\end{column}
\begin{column}{.32\textwidth}\centering
\textbf{Exp03 · + MixUp 0.2}\\\small Mantiene FT parcial\\\textbf{76,96\,\%}
\end{column}
\begin{column}{.32\textwidth}\centering
\textbf{Exp04 · FT completo}\\\small Conserva MixUp 0.2\\\textbf{77,97\,\%}: iguala scratch
\end{column}
\end{columns}
\end{frame}

% 9: 65 s
\begin{frame}{TL Challenge · Resolución, learning rate y MixUp}
\centering\small Se conservan FT completo salvo BN, batch 32 y WD $10^{-4}$
\grafica{tl_ajustes_finales.pdf}{3.95cm}
\begin{columns}[T,onlytextwidth]
\begin{column}{.32\textwidth}\centering
\textbf{Exp05 · 192 px}\\\small LR $10^{-5}$ · MixUp 0.2\\\textbf{80,16\,\%}: mejora
\end{column}
\begin{column}{.32\textwidth}\centering
\textbf{Exp06 · LR $3\cdot10^{-5}$}\\\small Mantiene 192 px y MixUp 0.2\\\textbf{80,91\,\%}: mejor accuracy
\end{column}
\begin{column}{.32\textwidth}\centering
\textbf{Exp07 · MixUp 0.4}\\\small Mantiene 192 px y LR\\\textbf{80,80\,\%}: menor CE
\end{column}
\end{columns}
\end{frame}

% 10: 45 s
\begin{frame}{TL Challenge · Errores del mejor modelo por accuracy}
\grafica{confusion_tl06.pdf}{5.25cm}
\end{frame}

% 11: 55 s
\begin{frame}{Balance · Mejores modelos de cada fase}
\grafica{balance_final.pdf}{4.6cm}
\idea{B0 TL exp06: \textbf{+2,93 pp} frente a B0 desde cero}
\end{frame}

% 12: 45 s
\begin{frame}{Cierre · Qué configuraciones conservamos}
\begin{center}
\begin{minipage}{.88\textwidth}
\begin{columns}[T,onlytextwidth]
\begin{column}{.47\textwidth}
\begin{block}{CNN Challenge}
\textbf{Popular: ResNet18 exp04}\\
192 px · smoothing 0.1\\[8pt]
\textbf{Custom: exp03}\\
128 px · MixUp 0.2
\end{block}
\end{column}
\begin{column}{.47\textwidth}
\begin{block}{TL Challenge · Candidatos a envío}
\textbf{B0 exp06: 80,91\,\%}\\
Mejor accuracy y F1\\[8pt]
\textbf{B0 exp07: 80,80\,\%}\\
Accuracy similar, menor CE
\end{block}
\end{column}
\end{columns}
\end{minipage}
\end{center}
\vspace{.3cm}
\idea{Conservar lo que mejora; descartar lo que no aporta}
\end{frame}
\end{document}
